% titel: Blatt 9, A5: Gerade und ungerade Funktionen
% skript: Kap. 5.2 (Hauptsatz, Stammfunktionen) und Kap. 5.3 (Substitution, S. 91)

% Teil I: Werkzeuge, kein Lösungsweg.

Symmetrie-Eigenschaften stetiger Funktionen und Konsequenzen für deren
Integrale über $[-a,a]$ sowie für die Parität von Stammfunktionen.

\begin{werkzeugkasten}[Gerade / ungerade]
$f:\mathbb R\to\mathbb R$ heißt \textbf{gerade}, falls $f(-x)=f(x)$, und
\textbf{ungerade}, falls $f(-x)=-f(x)$ für alle $x$.
\begin{itemize}
  \item Gerade: Graph symmetrisch zur $y$-Achse, $f$ durch Werte auf
        $[0,\infty)$ festgelegt.
  \item Ungerade: Graph punktsymmetrisch zum Ursprung.
  \item $f'\ $ dergerade, wenn $f$ ungerade ($f=\sin$); $f'$ ungerade,
        wenn $f$ gerade ($f=\cos$).
\end{itemize}
\textit{Skript:} Symmetrie-Beispiele, Kap. 2.6.
\end{werkzeugkasten}

\begin{werkzeugkasten}[Symmetrie im Integral]
Mit der Substitution $u=-x$ (Kap. 5.3) zerlegt man das Integral über
$[-a,a]$:
\begin{equation}
  \int_{-a}^{a} f(x)\,\dd x = \int_0^a f(x)\,\dd x + \int_{-a}^{0} f(x)\,\dd x .
\end{equation}
Für \textbf{gerade} $f$ fallen beide Teilflächen zusammen, für \textbf{ungerade}
heben sie sich auf. Dadurch lassen sich die Integrale vereinfachen.

\textit{Skript:} Substitutionsregel, Beispiel 5.3.4; Symmetrie in
Beispiel 5.2.3, S. 87.
\end{werkzeugkasten}

\begin{achtung}[Darauf achten]
\begin{itemize}
  \item In Teil c) die Stammfunktion gezielt wählen (hier $F(0)=0$) und die
        Parität von $F$ aus $F'=f$ plus Symmetrie ableiten.
  \item $f$ gerade $\Rightarrow\int_{-a}^{0}f(x)\dd x=\int_0^a f(x)\dd x$
        über $u=-x$; $f$ ungerade $\Rightarrow\int_{-a}^{0}f(x)\dd x=
        -\int_0^a f(x)\dd x$.
  \item Im Integral die Integrationsvariable umbenennen, damit alle
        Teilintegrale gleich heißen (z. B. nur noch $f(x)$), sonst vergleicht
        man unpassende Ausdrücke.
  \item Die Behauptung in c) ist die \emph{Umkehrung} des bekannten
        Zusammenhangs für Ableitungen — nicht verwechseln.
\end{itemize}
\end{achtung}
